\chapter{Celiac Disease}
\index{Celiac Disease}
\label{sec:Celiac Disease}

\section{Overview}
Celiac disease is an autoimmune enteropathy triggered by ingestion of gluten (gliadin and glutenin proteins found in wheat, barley, and rye) in genetically susceptible individuals (HLA-DQ2 or HLA-DQ8). Prevalence is approximately 1\% worldwide.
\section{Causes}
This condition happens because of an immune response causing villous shrinkage or wasting, crypt abnormal increase in cells, and intraepithelial lymphocytosis in the small intestinal mucosa, primarily affecting the duodenum and jejunum. The underlying mechanisms involve complex interactions between genetic predisposition and environmental factors that disrupt normal physiological processes. The underlying mechanisms involve complex interactions between genetic predisposition and environmental factors. Disruption of normal physiological processes leads to the characteristic features of the condition. Multiple contributing factors may act synergistically to trigger or worsen the condition. Understanding the specific cause helps guide targeted treatment approaches.
\section{Risk Factors}


\begin{itemize}[noitemsep]
  \item Genetic predisposition and family history
  \item Advancing age
  \item Lifestyle factors (diet, exercise, smoking, alcohol)
  \item Environmental exposures
  \item Underlying medical conditions
  \item Medication use
\end{itemize}
\section{Signs and Symptoms}

\subsection{Common symptoms}
\begin{itemize}[noitemsep]
  \item Clinical features vary widely: classical symptoms in children include chronic diarrhea, failure to thrive, and abdominal distension
  \item In adults, diarrhea, weight loss, iron deficiency anemia, and osteoporosis. Dermatitis herpetiformis is the cutaneous manifestation. Many patients are asymptomatic or have extraintestinal manifestations (fatigue, neurological symptoms, liver disease, infertility).
  \item Pain or discomfort in the affected area
  \end{itemize}

\subsection{Emergency warning signs}
\begin{itemize}[noitemsep]
  \item Severe or worsening symptoms of celiac disease require immediate medical evaluation
\end{itemize}
\section{Progression and Stages}
The condition may progress through different stages of severity over time. Early stages often involve mild or subtle symptoms that may be overlooked. Moderate stages involve more noticeable symptoms affecting daily function. Advanced stages may involve significant impairment and complications.
\section{Complications}
\begin{itemize}[noitemsep]
  \item Disease progression leading to worsening symptoms
  \item Secondary infections or injuries
  \item Side effects of treatment
  \item Reduced quality of life
  \item Emotional and psychological impact
\end{itemize}
\section{Diagnosis}
\begin{itemize}[noitemsep]
  \item To make the diagnosis, doctors look for serological testing (anti-tissue transglutaminase IgA, anti-endomysial antibodies) and duodenal biopsy showing villous shrinkage or wasting, crypt abnormal increase in cells, and increased intraepithelial lymphocytes
  \item HLA testing helps exclude the diagnosis.
  \item Comprehensive medical history and physical examination
  \item Laboratory tests to assess organ function
  \item Imaging studies to visualize affected structures
  \item Specialized testing depending on the affected system
  \item Consultation with relevant specialists
\end{itemize}
\section{Prevention}
\begin{itemize}[noitemsep]
  \item Maintain a healthy lifestyle with balanced diet and exercise
  \item Avoid known risk factors
  \item Attend regular medical check-ups for early detection
  \item Follow recommended screening guidelines
\end{itemize}
\section{Treatment Options}
Treatment typically involves a strict lifelong gluten-free diet, which leads to mucosal recovery and symptom resolution. Medications to manage symptoms and address underlying mechanisms Lifestyle modifications and self-care measures Physical therapy and rehabilitation Surgical intervention when indicated Regular monitoring and follow-up care Patient education and support services

\begin{itemize}[noitemsep]
  \item Medications to manage symptoms and address underlying mechanisms
  \item Lifestyle modifications and self-care measures
  \item Physical therapy and rehabilitation
  \item Surgical intervention when indicated
  \item Regular monitoring and follow-up care
\end{itemize}
\section{Living with It}
\begin{itemize}[noitemsep]
  \item Follow your treatment plan as prescribed
  \item Maintain regular follow-up with your healthcare provider
  \item Make necessary lifestyle adjustments
  \item Seek support from family, friends, or support groups
  \item Stay informed about your condition
\end{itemize}
\section{Outlook}
Most people recover well with dietary compliance. The outlook varies depending on the specific condition, severity at diagnosis, and response to treatment. Early intervention generally leads to better outcomes. Many conditions can be effectively managed with appropriate treatment. Regular follow-up care is important for monitoring and treatment adjustment.
